\subsection{21.121(b): sharp soluble bounds and permanence}
\label{partial:jordan}
\problemstatement{21.121}

Use the definition of admissible \(p\)-Jordan exponents from
Section~\ref{cand:21.121}. The unresolved part asks for a universal
upper bound three. The following arguments prove exponent two
under a soluble-subgroup hypothesis, give a nonsoluble example
with exact exponent strictly between two and three, and preserve
the entire set of admissible exponents under specified operations.
The definition and prior general framework are those in
\cite{hu-jordan,chen-shramov}.

\paragraph{A characteristic abelian subgroup lemma.}
If a finite group \(X\) has an abelian subgroup of index at most
\(j\), it has a characteristic abelian subgroup of index at most
\(j^2\). Recall the Chermak--Delgado argument in this form.
For \(H\le X\) put \(m(H)=|H||C_X(H)|\).
The product-size formula and centralizer inclusions give
\[
 m(H)m(K)\le m(H\cap K)m(\langle H,K\rangle).
\]
Thus the intersection \(D\) of all subgroups of maximum measure
also has maximum measure and is characteristic. Its centralizer
has maximum measure as well, since \(D\subseteq C_X(C_X(D))\).
Minimality of the intersection implies \(D\subseteq C_X(D)\),
so \(D\) is abelian. For the given abelian \(A\),
\[
 |X||D|\ge m(D)\ge m(A)\ge |A|^2\ge |X|^2/j^2,
\]
proving the bound. The lemma does not force \(D\subseteq A\).
It will be applied inside a \(p'\)-group when prime-to-\(p\)
order is required.

An abelian \(p'\)-subgroup \(C\le\Aut(P)\), for a nontrivial
finite \(p\)-group \(P\), satisfies
\[
 |C|\le |P/\Phi(P)|-1.
\]
Indeed the kernel of \(\Aut(P)\to\Aut(P/\Phi(P))\) is a \(p\)-group:
it acts freely on the \(|\Phi(P)|^d\) tuples lifting a fixed
basis of \(P/\Phi(P)\), because each lift tuple generates \(P\).
Thus \(C\) acts faithfully on \(\F_p^d\). Complete reducibility
and commutativity embed it in a product of multiplicative groups
of finite fields of orders \(p^{d_i}\), with \(\sum d_i=d\).
Their orders have product at most \(p^d-1\).

\paragraph{The soluble theorem.}
Suppose every finite \(p'\)-subgroup of \(\Gamma\) has a normal
abelian subgroup of index at most \(J\), and every finite
subgroup of \(\Gamma\) is soluble. For each finite \(G\le\Gamma\)
there is a characteristic abelian \(p'\)-subgroup \(B\) with
\[
 [G:B]\le J^3|G_{(p)}|^2.
\]
First let \(R\ne1\) be finite soluble with \(O_{p'}(R)=1\).
Its Fitting subgroup \(F(R)\ne1\) is a \(p\)-group and
\(C_R(F(R))\subseteq F(R)\).
For the latter standard fact, if the centralizer modulo \(Z(F(R))\)
were nontrivial, choose a minimal normal elementary abelian
subgroup there. Its preimage is a normal class-two nilpotent
subgroup strictly exceeding the centralizer's intersection with
\(F(R)\), contradicting the Fitting theorem.
A Hall \(p'\)-subgroup \(H\) therefore acts faithfully on \(F(R)\).
Choose \(A\lhd H\) abelian of index at most \(J\) and apply the
coprime-action bound:
\[
 |R|=|R_{(p)}||H|
 \le J|R_{(p)}|(|R_{(p)}|-1).
\]
Hall-subgroup existence here is the soluble Hall theorem,
obtainable by induction with Schur--Zassenhaus.

For general \(G\), set \(N=O_{p'}(G)\) and \(R=G/N\).
Then \(O_{p'}(R)=1\); the same \(J\) works for its \(p'\)-subgroups,
whose full preimages are \(p'\)-groups. The characteristic lemma
gives \(B\mathrel{\mathrm{char}} N\) abelian with \([N:B]\le J^2\).
Since \(N\mathrel{\mathrm{char}} G\), also \(B\mathrel{\mathrm{char}} G\), and the preceding estimate
proves the assertion, with \(R=1\) handled directly.
If every finite subgroup instead has a soluble subgroup of
index at most \(d\), take its core, of index at most \(d!\).
The bound becomes \(d!J^3|G_{(p)}|^2\), with \(B\) normal in \(G\).

When every finite subgroup has normal Sylow \(p\)-subgroup \(P\),
a sharper bound is available without solubility.
Write \(G=P\rtimes H\) by Schur--Zassenhaus, choose
\(A\lhd H\) abelian with \([H:A]\le J\), and put \(B=C_A(P)\).
Both \(P\) and \(H\) normalize \(B\), so \(B\lhd G\).
For \(P\ne1\),
\[
 [G:B]\le J|P|(|P/\Phi(P)|-1)<J|P|^2;
\]
for \(P=1\), use \(A\). Exponent two is sharp for every \(p\):
the locally finite affine group \(\operatorname{AGL}_1(\overline{\F}_p)\)
has these hypotheses with \(J=1\). Its subgroups
\(\operatorname{AGL}_1(\F_q)\), \(q=p^n\), have order \(q(q-1)\)
and no nontrivial normal \(p'\)-subgroup.
Such a subgroup would centralize the normal translations, whose
centralizer is the translation group itself.
Hence their least normal abelian \(p'\)-indices are \(q(q-1)\),
excluding every exponent less than two.

\paragraph{An exact nonsoluble exponent.}
Let \(\Gamma=\bigoplus_{\mathbb N}A_5\), the finite-support direct
product, and \(\rho=\log(60)/\log(4)\).
Then its \(2\)-Jordan exponent is exactly \(\rho\), attained with
constant one. For a subgroup \(H\le A_5^n\), separate coordinates
with full projection from those with proper projection.
The first projection \(T\) is a product of diagonal simple
factors, hence \(A_5^r\); induction proves this by projecting off
the last coordinate, whose kernel is either \(A_5\) or trivial.
In the latter case a surjection \(A_5^r\to A_5\) is supported
on exactly one factor because its factor images commute and
are normal. The second projection \(B\) is soluble, since
all proper subgroups of \(A_5\) are soluble.
Goursat's argument gives a common quotient of \(T\) and \(B\);
it is perfect and soluble, hence trivial. Thus \(H=T\times B\).
Every odd-order subgroup of \(B\) is abelian, since its coordinate
projections are cyclic subgroups of orders \(1,3,5\).
The soluble theorem with \(J=1\) gives a normal abelian odd-order
\(D\le B\) of index at most \(|B_{(2)}|^2\). Therefore
\[
 [H:1\times D]\le60^r|B_{(2)}|^2
 \le(4^r|B_{(2)}|)^\rho=|H_{(2)}|^\rho.
\]
Every finite subgroup of \(\Gamma\) has finite support, so this
is a uniform bound. Full powers \(A_5^r\) have no nontrivial
abelian normal subgroup and exclude every smaller exponent.
Since \(16<60<64\), one has \(2<\rho<3\).
This example is not virtually soluble: a soluble subgroup of
\(A_5^r\) has index at least \(5^r\).

\paragraph{Preserving exponents.}
Suppose \(e\) is admissible for \(\Gamma\) with constant \(J\).
For any finite \(G\le\Gamma\), apply the characteristic lemma
inside \(N=O_{p'}(G)\). The original abelian normal subgroup
lies in \(N\), giving \([G:N]\le J|G_{(p)}|^e\), and the
\(p'\)-subgroup bound gives a characteristic abelian \(D\le N\)
of index at most \(J^2\). Thus
\[
 [G:D]\le J^3|G_{(p)}|^e.
\]
Replacing normal by characteristic in the definition does not
change its admissible exponents or attainment.
If \(\Lambda\lhd\Gamma\) and finite subgroups of
\(\Gamma/\Lambda\) have order at most \(b\), apply this form
to \(G\cap\Lambda\); its characteristic subgroup is normal in
\(G\). For a real \(e\), the constant
\(bJ^3\max(1,b^{-e})\) works in \(\Gamma\).
Restriction gives the reverse implication. Taking a core proves
the same invariance for any finite-index subgroup.

More generally, assume a fixed generalized \(p\)-Jordan constant
\(J_0\) for \(\Gamma\), and let \(N\lhd\Gamma\) be finite.
Then \(\Gamma\) and \(\Gamma/N\) have exactly the same admissible
exponents. Downward passage uses finite full preimages.
For upward passage, take finite \(G\le\Gamma\), \(K=G\cap N\),
and the preimage \(M\) of a normal abelian \(p'\)-subgroup
in \(G/K\). Put \(C=C_M(K)\). It is normal in \(G\),
\([M:C]\le|N|!\), and \(C'\subseteq K\subseteq C_G(C)\);
so \(C\) has class at most two.
The soluble theorem supplies \(D\mathrel{\mathrm{char}} C\), abelian of \(p'\)-order,
with \([C:D]\le J_0^3|K_{(p)}|^2\).
If the quotient bound has constant \(J\), then
\[
 [G:D]\le
 JJ_0^3|N|!\max(1,|N|^{2-e})\,|G_{(p)}|^e.
\]
This preserves the exponent and its attainment.
For a finite normal \(p\)-kernel the generalized hypothesis
follows from the quotient bound, since \(p'\)-subgroups inject
into the quotient. For other finite kernels it is necessary:
the infinite extraspecial group at a prime \(q\ne p\) has
central kernel \(C_q\) and abelian quotient, but contains finite
extraspecial groups with unbounded least abelian index.

These results leave the unrestricted bound three unresolved.
The full arguments and constants are in
\artifact{research/21.121b-soluble-proof.md},
\artifact{research/21.121b-normal-sylow-proof.md},
\artifact{research/21.121b-a5-proof.md},
\artifact{research/21.121-permanence-proof.md}, and
\artifact{research/21.121-finite-kernels.md}.
